\chapter{Background}%
\label{ch:background}

In this chapter, you provide information that has a benefit for an informed audience.
Often, you use this to explain a technique or a concept you use, referring to the literature.
This is different from Related Work (\cref{ch:related}), where you show the relation of literature to your contributions (to ground, compare, or inform).
A good rule of thumb is that your thesis should still make sense if the background chapter is removed.

\begin{figure}[htbp]
\centering
\resizebox{\textwidth}{!}{%
\begin{tikzpicture}[
    thesisfig,
    node distance=0.30cm,
    image/.style={tfbox},
    stem/.style ={tfbox, draw=orange!75!black, fill=orange!15},
    layer/.style={tfbox, draw=blue!65!black,   fill=blue!14},
    head/.style ={tfbox, draw=green!45!black,  fill=green!12},
    arrow/.style={tfarrow}
]

%% top-level ResNet-18 pipeline: stem and head collapsed to single boxes;
%% their internals (Conv-BN-ReLU-MaxPool / AvgPool-FC) are given in the caption.
\node[image] (img) {Image};
\node[stem, right=of img] (stem) {Stem};

\node[layer, right=0.5cm of stem] (l1) {Layer 1};
\node[layer, right=of l1] (l2) {Layer 2};
\node[layer, right=of l2] (l3) {Layer 3};
\node[layer, right=of l3] (l4) {Layer 4};

\node[head, right=0.5cm of l4] (head) {Head};

\draw[arrow] (img) -- (stem);
\draw[arrow] (stem) -- (l1);
\draw[arrow] (l1) -- (l2);
\draw[arrow] (l2) -- (l3);
\draw[arrow] (l3) -- (l4);
\draw[arrow] (l4) -- (head);

\draw[tfbrace]
    ($(l1.north west)+(0,0.12)$) -- ($(l4.north east)+(0,0.12)$)
    node[midway, yshift=0.45cm] {Residual stages};

\end{tikzpicture}%
%
}
\caption{Top-level ResNet-18 inference pipeline used in this thesis. The model is treated as a stem, four residual stages, and a classification head.}
\label{fig:resnet18_pipeline}
\end{figure}

ResNet-18 can be viewed as three major parts for the purposes of this thesis:
an initial stem, four residual stages, and a classification head. The stem performs
the first feature extraction and spatial downsampling through convolution,
normalisation, activation, and max pooling. The four residual stages, denoted
\texttt{layer1} through \texttt{layer4}, contain the main residual-block computation
of the network. The classification head then reduces the final feature map through
average pooling and applies the final fully connected classifier.

This structure motivates the coarse split choices used later in RQ1.1. The thesis
does not treat every primitive operation as an equally meaningful split candidate.
Instead, the initial coarse sweep uses the four residual-stage boundaries after
\texttt{layer1}, \texttt{layer2}, \texttt{layer3}, and \texttt{layer4}. These boundaries
preserve the architectural stage structure of ResNet-18 and provide a controlled
progression from early, high-activation regions to later, lower-activation regions.
Splits inside the stem are excluded from the primary coarse sweep because they
would place boundaries between tightly coupled primitive operations before
substantial activation-size reduction. Splits inside the classification head are also
excluded because the head contains only the final pooling and linear classification
operations, leaving one side of the partition with negligible computation.

%%split point image
\begin{figure}[htbp]
\centering
\resizebox{\textwidth}{!}{%
\begin{tikzpicture}[
    thesisfig,
    node distance=0.45cm,
    block/.style={tfbox, minimum width=1.0cm, inner xsep=3pt},
    stem/.style ={tfbox, minimum width=1.0cm, inner xsep=3pt, draw=orange!75!black, fill=orange!15},
    layer/.style={tfbox, minimum width=1.0cm, inner xsep=3pt, draw=blue!65!black,   fill=blue!14},
    head/.style ={tfbox, minimum width=1.0cm, inner xsep=3pt, draw=green!45!black,  fill=green!12},
    arrow/.style={tfarrow},
    split/.style={dashed, very thick, draw=red!75!black},
    splitlabel/.style={font=\small, text=red!75!black, align=center}
]

% Main pipeline
\node[block] (img) {Image};
\node[stem, right=of img] (stem) {Stem};
\node[layer, right=of stem] (l1) {Layer 1};
\node[layer, right=of l1] (l2) {Layer 2};
\node[layer, right=of l2] (l3) {Layer 3};
\node[layer, right=of l3] (l4) {Layer 4};
\node[head, right=of l4] (head) {Classification\\Head};

% Arrows
\draw[arrow] (img) -- (stem);
\draw[arrow] (stem) -- (l1);
\draw[arrow] (l1) -- (l2);
\draw[arrow] (l2) -- (l3);
\draw[arrow] (l3) -- (l4);
\draw[arrow] (l4) -- (head);

% Stage~L1 coarse split markers after each residual stage
\draw[split] ($(l1.east)+(0.15,0.75)$) -- ($(l1.east)+(0.15,-0.75)$);
\draw[split] ($(l2.east)+(0.15,0.75)$) -- ($(l2.east)+(0.15,-0.75)$);
\draw[split] ($(l3.east)+(0.15,0.75)$) -- ($(l3.east)+(0.15,-0.75)$);
\draw[split] ($(l4.east)+(0.15,0.75)$) -- ($(l4.east)+(0.15,-0.75)$);

% Labels under split markers
\node[splitlabel, below=0.75cm of l1.east, xshift=0.15cm] {split after\\layer1};
\node[splitlabel, below=0.75cm of l2.east, xshift=0.15cm] {split after\\layer2};
\node[splitlabel, below=0.75cm of l3.east, xshift=0.15cm] {split after\\layer3};
\node[splitlabel, below=0.75cm of l4.east, xshift=0.15cm] {split after\\layer4};

% Optional RQ label
\node[
    font=\normalsize\bfseries,
    text=red!75!black,
    above=0.65cm of l2,
    xshift=0.95cm
] {Stage~L1 coarse split candidates};

\end{tikzpicture}%%
}
\caption{Coarse ResNet-18 partition boundaries evaluated in RQ1.1. The stem and classification head are collapsed into single blocks, while the four residual stages are shown as the candidate split locations.}
\label{fig:rq11_coarse_split_boundaries}
\end{figure}

As shown in \cref{fig:rq11_coarse_split_boundaries}, RQ1.1 evaluates split
points after the four major residual stages rather than inside the stem or
classification head.

%%segmented image of resnet18 after layer2 split
\begin{figure}[htbp]
\centering
\resizebox{\textwidth}{!}{%
\begin{tikzpicture}[
    thesisfig,
    node distance=0.24cm,
    block/.style={tfbox, minimum width=1.2cm},
    stem/.style ={tfbox, minimum width=1.2cm, draw=orange!75!black, fill=orange!15},
    layer/.style={tfbox, minimum width=1.2cm, draw=blue!65!black,   fill=blue!14},
    head/.style ={tfbox, minimum width=1.2cm, draw=green!45!black,  fill=green!12},
    arrow/.style={tfarrow},
    split/.style={dashed, very thick, draw=red!75!black},
    segmentA/.style={draw=orange!75!black, fill=orange!7, rounded corners=5pt, thick, inner sep=5pt},
    segmentB/.style={draw=blue!65!black,   fill=blue!6,  rounded corners=5pt, thick, inner sep=5pt}
]

% Main pipeline
\node[block] (img) {Image};
\node[stem, right=of img] (stem) {Stem};
\node[layer, right=of stem] (l1) {Layer 1};
\node[layer, right=of l1] (l2) {Layer 2};

% Larger gap around the split
\node[layer, right=0.7cm of l2] (l3) {Layer 3};
\node[layer, right=of l3] (l4) {Layer 4};
\node[head, right=of l4] (head) {Classification\\Head};

% Split coordinate halfway between Layer 2 and Layer 3
\coordinate (splitline) at ($(l2.east)!0.5!(l3.west)$);

% Segment backgrounds
\begin{pgfonlayer}{background}
\node[
    segmentA,
    fit=(img)(stem)(l1)(l2),
    label={[font=\normalsize, text=orange!75!black]above:Segment 1 / Service A}
] (segA) {};

\node[
    segmentB,
    fit=(l3)(l4)(head),
    label={[font=\normalsize, text=blue!65!black]above:Segment 2 / Service B}
] (segB) {};
\end{pgfonlayer}

% Arrows
\draw[arrow] (img) -- (stem);
\draw[arrow] (stem) -- (l1);
\draw[arrow] (l1) -- (l2);

\draw[arrow]
    (l2) -- node[
        midway,
        above=1.05cm of splitline,
        font=\normalsize,
        text=red!75!black
    ] {activation tensor}
    (l3);

\draw[arrow] (l3) -- (l4);
\draw[arrow] (l4) -- (head);

% Split marker after layer2
\draw[split]
    ($(splitline)+(0,1)$) -- ($(splitline)+(0,-1)$);

\node[
    font=\normalsize,
    text=red!75!black,
    align=center,
    below=1.05cm of splitline
] {split point\\after layer2};

\end{tikzpicture}%%
}
\caption{Example two-segment decomposition produced by a split after \texttt{layer2}. The first segment executes the stem through \texttt{layer2}; the second segment receives the intermediate activation tensor and executes \texttt{layer3} through the classification head.}
\label{fig:split_after_layer2_two_segments}
\end{figure}
A split boundary creates two execution segments, as illustrated for
\texttt{split\_after\_layer2} in \cref{fig:split_after_layer2_two_segments}.

%% service-boundary / RPC cost figure
\begin{figure}[htbp]
\centering
\resizebox{0.96\textwidth}{!}{%
\begin{tikzpicture}[
    font=\small,
    service/.style={
        draw=black!60,
        fill=black!3,
        rounded corners=6pt,
        thick,
        inner xsep=16pt,
        inner ysep=14pt
    },
    stepA/.style={
        draw=orange!75!black,
        fill=orange!14,
        rounded corners=3pt,
        minimum width=2.6cm,
        minimum height=0.85cm,
        align=center,
        thick
    },
    stepB/.style={
        draw=blue!65!black,
        fill=blue!12,
        rounded corners=3pt,
        minimum width=2.6cm,
        minimum height=0.85cm,
        align=center,
        thick
    },
    rpc/.style={
        draw=red!70!black,
        fill=red!8,
        rounded corners=3pt,
        minimum width=2.9cm,
        minimum height=0.85cm,
        align=center,
        thick
    },
    arrow/.style={
        -{Latex[length=2mm]},
        thick,
        draw=black!70
    },
    returnarrow/.style={
        -{Latex[length=2mm]},
        thick,
        draw=black!55,
        dashed
    },
    boundary/.style={
        dashed,
        very thick,
        draw=red!75!black
    },
    labeltext/.style={
        font=\scriptsize,
        text=black!65,
        align=center
    }
]

% Top row: forward path
\node[stepA] (a_compute) {compute\\Segment 1};
\node[stepA, right=0.75cm of a_compute] (a_ser)    {serialize\\activation};

\node[rpc,   right=1.5cm  of a_ser]     (grpc_req) {gRPC request\\activation tensor};

\node[stepB, right=1.5cm  of grpc_req]  (b_deser)  {deserialize\\activation};
\node[stepB, right=0.75cm of b_deser]   (b_compute){compute\\Segment 2};
\node[stepB, right=0.75cm of b_compute] (b_ser)    {serialize\\output};

% Bottom row: return path
% a_deser sits directly below a_ser so the columns stay aligned.
% grpc_resp sits directly below grpc_req so the boundary stays vertical.
\node[stepA, below=2.5cm of a_ser]     (a_deser)   {deserialize\\output};
\node[rpc,   below=2.5cm of grpc_req]  (grpc_resp) {gRPC response\\output tensor};

% Service background boxes
\begin{pgfonlayer}{background}
\node[
    service,
    fit=(a_compute)(a_ser)(a_deser),
    label={[font=\small\bfseries, text=orange!75!black]above:Service A / Segment 1}
] (serviceA) {};

\node[
    service,
    fit=(b_deser)(b_compute)(b_ser),
    label={[font=\small\bfseries, text=blue!65!black]above:Service B / Segment 2}
] (serviceB) {};
\end{pgfonlayer}

% Forward arrows
\draw[arrow] (a_compute) -- (a_ser);
\draw[arrow] (a_ser)     -- (grpc_req);
\draw[arrow] (grpc_req)  -- (b_deser);
\draw[arrow] (b_deser)   -- (b_compute);
\draw[arrow] (b_compute) -- (b_ser);

% Return arrows
\draw[returnarrow] (b_ser.south)  |- (grpc_resp.east);
\draw[returnarrow] (grpc_resp)    -- (a_deser);

% Service boundary line
\coordinate (boundarytop)
    at ($(a_ser.east)!0.5!(grpc_req.west)+(0,1.1)$);
\coordinate (boundarybottom)
    at ($(a_deser.east)!0.5!(grpc_resp.west)+(0,-1.1)$);

\draw[boundary] (boundarytop) -- (boundarybottom);

\node[
    font=\scriptsize\bfseries,
    text=red!75!black,
    align=center,
    above=0.08cm of boundarytop
] {service\\boundary};

\end{tikzpicture}%%
}
\caption{Runtime cost components introduced by a two-service inference
boundary. In addition to the model computation on each side of the split,
the intermediate activation must be serialized, transferred through gRPC,
deserialized by the downstream service, and followed by response
serialization and reconstruction. The boundary-crossing interval therefore
includes serialization, RPC transport, remote execution, and response
reconstruction.}
\label{fig:service_boundary_rpc_cost}
\end{figure}

%%split inside layer3 image
\begin{figure}[htbp]
\centering
\resizebox{\textwidth}{!}{%
\begin{tikzpicture}[
    thesisfig,
    node distance=0.36cm,                 % gap *within* a layer's block pair
    image/.style={tfbox, minimum width=0.7cm, inner xsep=3pt},
    stem/.style ={tfbox, minimum width=0.7cm, inner xsep=3pt, draw=orange!75!black, fill=orange!15},
    block/.style={tfbox, minimum width=0.55cm, inner xsep=3pt, draw=blue!65!black,   fill=blue!14},
    head/.style ={tfbox, minimum width=0.7cm, inner xsep=3pt, draw=green!45!black,  fill=green!12},
    arrow/.style={tfarrow},
    layerbrace/.style={decorate, decoration={brace, amplitude=5pt}, thick, draw=black!65}
]

% Main pipeline. Blocks abbreviated B0/B1 (zero-indexed, as in layerN.0/.1;
% defined in caption) so all 11 boxes fit one row at body-text size; layers
% grouped by the braces above.
\node[image] (img) {Image};
\node[stem, right=0.38cm of img] (stem) {Stem};

% Layer 1 blocks
\node[block, right=0.45cm of stem] (l1b1) {B0};
\node[block, right=of l1b1] (l1b2) {B1};

% Layer 2 blocks
\node[block, right=0.45cm of l1b2] (l2b1) {B0};
\node[block, right=of l2b1] (l2b2) {B1};

% Layer 3 blocks
\node[block, right=0.45cm of l2b2] (l3b1) {B0};
\node[block, right=of l3b1] (l3b2) {B1};

% Layer 4 blocks
\node[block, right=0.45cm of l3b2] (l4b1) {B0};
\node[block, right=of l4b1] (l4b2) {B1};

\node[head, right=0.45cm of l4b2] (head) {Head};

% Arrows
\draw[arrow] (img) -- (stem);
\draw[arrow] (stem) -- (l1b1);
\draw[arrow] (l1b1) -- (l1b2);
\draw[arrow] (l1b2) -- (l2b1);
\draw[arrow] (l2b1) -- (l2b2);
\draw[arrow] (l2b2) -- (l3b1);
\draw[arrow] (l3b1) -- (l3b2);
\draw[arrow] (l3b2) -- (l4b1);
\draw[arrow] (l4b1) -- (l4b2);
\draw[arrow] (l4b2) -- (head);

% Layer labels using braces
\draw[layerbrace]
    ($(l1b1.north west)+(0,0.18)$) -- ($(l1b2.north east)+(0,0.18)$)
    node[midway, yshift=0.45cm] {Layer 1};

\draw[layerbrace]
    ($(l2b1.north west)+(0,0.18)$) -- ($(l2b2.north east)+(0,0.18)$)
    node[midway, yshift=0.45cm] {Layer 2};

\draw[layerbrace]
    ($(l3b1.north west)+(0,0.18)$) -- ($(l3b2.north east)+(0,0.18)$)
    node[midway, yshift=0.45cm] {Layer 3};

\draw[layerbrace]
    ($(l4b1.north west)+(0,0.18)$) -- ($(l4b2.north east)+(0,0.18)$)
    node[midway, yshift=0.45cm] {Layer 4};

\end{tikzpicture}%
%
}
\caption{Simplified ResNet-18 block structure. The stem and classification head are shown as single components, while each residual stage is shown as two basic blocks.}
\label{fig:resnet18_block_structure}
\end{figure}

%%split inside layer3 image with split after block0
\begin{figure}[htbp]
\centering
\resizebox{\textwidth}{!}{%
\begin{tikzpicture}[
    font=\small,
    image/.style={
        draw=black!55,
        fill=black!5,
        rounded corners=3pt,
        minimum width=1.25cm,
        minimum height=0.85cm,
        align=center,
        thick
    },
    stem/.style={
        draw=orange!75!black,
        fill=orange!15,
        rounded corners=3pt,
        minimum width=1.35cm,
        minimum height=0.85cm,
        align=center,
        thick
    },
    block/.style={
        draw=blue!65!black,
        fill=blue!14,
        rounded corners=3pt,
        minimum width=1.25cm,
        minimum height=0.85cm,
        align=center,
        thick
    },
    head/.style={
        draw=green!45!black,
        fill=green!12,
        rounded corners=3pt,
        minimum width=1.65cm,
        minimum height=0.85cm,
        align=center,
        thick
    },
    arrow/.style={
        -{Latex[length=2mm]},
        thick,
        draw=black!70
    },
    layerbrace/.style={
        decorate,
        decoration={brace, amplitude=5pt},
        thick,
        draw=black!65
    },
    split/.style={
        dashed,
        very thick,
        draw=red!75!black
    },
    splitlabel/.style={
        font=\scriptsize\bfseries,
        text=red!75!black,
        align=center
    },
    segmentA/.style={
        draw=orange!75!black,
        fill=orange!7,
        rounded corners=6pt,
        thick,
        inner xsep=12pt,
        inner ysep=28pt
    },
    segmentB/.style={
        draw=blue!65!black,
        fill=blue!6,
        rounded corners=6pt,
        thick,
        inner xsep=12pt,
        inner ysep=28pt
    },
    node distance=0.35cm
]

% Main pipeline
\node[image] (img) {Image};
\node[stem, right=of img] (stem) {Stem};

% Layer 1 blocks
\node[block, right=0.65cm of stem] (l1b1) {Block 1};
\node[block, right=of l1b1] (l1b2) {Block 2};

% Layer 2 blocks
\node[block, right=0.65cm of l1b2] (l2b1) {Block 1};
\node[block, right=of l2b1] (l2b2) {Block 2};

% Layer 3 blocks, with extra spacing around the internal split
\node[block, right=0.80cm of l2b2] (l3b1) {Block 1};
\node[block, right=2.75cm of l3b1] (l3b2) {Block 2};

% Layer 4 blocks
\node[block, right=0.65cm of l3b2] (l4b1) {Block 1};
\node[block, right=of l4b1] (l4b2) {Block 2};

\node[head, right=0.65cm of l4b2] (head) {Classification\\Head};

% Split coordinate inside Layer 3
\coordinate (splitline) at ($(l3b1.east)!0.5!(l3b2.west)$);

% Segment backgrounds
\begin{pgfonlayer}{background}
\node[
    segmentA,
    fit=(img)(stem)(l1b1)(l1b2)(l2b1)(l2b2)(l3b1)
] (segA) {};

\node[
    segmentB,
    fit=(l3b2)(l4b1)(l4b2)(head)
] (segB) {};
\end{pgfonlayer}

% Segment labels
\node[
    font=\small\bfseries,
    text=orange!75!black,
    above=0.38cm of segA.north
] {Segment 1 / Service A};

\node[
    font=\small\bfseries,
    text=blue!65!black,
    above=0.38cm of segB.north
] {Segment 2 / Service B};

% Arrows
\draw[arrow] (img) -- (stem);
\draw[arrow] (stem) -- (l1b1);
\draw[arrow] (l1b1) -- (l1b2);
\draw[arrow] (l1b2) -- (l2b1);
\draw[arrow] (l2b1) -- (l2b2);
\draw[arrow] (l2b2) -- (l3b1);
\draw[arrow] (l3b1) -- (l3b2);
\draw[arrow] (l3b2) -- (l4b1);
\draw[arrow] (l4b1) -- (l4b2);
\draw[arrow] (l4b2) -- (head);

% Layer labels using braces, placed inside the segment boxes
\draw[layerbrace]
    ($(l1b1.north west)+(0,0.10)$) -- ($(l1b2.north east)+(0,0.10)$)
    node[midway, yshift=0.36cm] {\small Layer 1};

\draw[layerbrace]
    ($(l2b1.north west)+(0,0.10)$) -- ($(l2b2.north east)+(0,0.10)$)
    node[midway, yshift=0.36cm] {\small Layer 2};

\draw[layerbrace]
    ($(l3b1.north west)+(0,0.10)$) -- ($(l3b2.north east)+(0,0.10)$)
    node[midway, yshift=0.36cm] {\small Layer 3};

\draw[layerbrace]
    ($(l4b1.north west)+(0,0.10)$) -- ($(l4b2.north east)+(0,0.10)$)
    node[midway, yshift=0.36cm] {\small Layer 4};

% Split marker after Layer 3 Block 1
\draw[split]
    ($(splitline)+(0,0.5)$) -- ($(splitline)+(0,-0.5)$);

\node[
    splitlabel,
    below=0.50cm of splitline
] {Stage~L2 split point\\after \texttt{layer3.0}};

\end{tikzpicture}%%
}
\caption{Block-level two-segment decomposition used in RQ1.2. The refined split point is placed inside \texttt{layer3}, after the first residual block. Segment 1 executes the stem through \texttt{layer3.0}; Segment 2 receives the intermediate activation tensor and executes \texttt{layer3.1} through the classification head.}
\label{fig:rq12_split_after_layer3_block0_segments}
\end{figure}

%%all detailed layers 
\begin{figure}[htbp]
\centering
\resizebox{\textwidth}{!}{%
\begin{tikzpicture}[     font=\normalsize,
   stem/.style={
        draw=black!65,
        fill=red!18,
        rounded corners=2pt,
        minimum width=0.72cm,
        minimum height=2.6cm,
        align=center,
        thick
    },
    pool/.style={
        draw=black!65,
        fill=black!8,
        rounded corners=2pt,
        minimum width=0.72cm,
        minimum height=2.6cm,
        align=center,
        thick
    },
    lone/.style={
        draw=cyan!60!black,
        fill=cyan!20,
        rounded corners=2pt,
        minimum width=0.72cm,
        minimum height=2.6cm,
        align=center,
        thick
    },
    ltwo/.style={
        draw=orange!65!black,
        fill=orange!18,
        rounded corners=2pt,
        minimum width=0.72cm,
        minimum height=2.6cm,
        align=center,
        thick
    },
    lthree/.style={
        draw=yellow!55!black,
        fill=yellow!25,
        rounded corners=2pt,
        minimum width=0.72cm,
        minimum height=2.6cm,
        align=center,
        thick
    },
    lfour/.style={
        draw=green!45!black,
        fill=green!18,
        rounded corners=2pt,
        minimum width=0.72cm,
        minimum height=2.6cm,
        align=center,
        thick
    },
    head/.style={
        draw=black!65,
        fill=black!12,
        rounded corners=2pt,
        minimum width=0.78cm,
        minimum height=2.6cm,
        align=center,
        thick
    },
    basicblockbox/.style={
        draw=black!45,
        fill=white,
        rounded corners=5pt,
        thick,
        inner xsep=10pt,
        inner ysep=8pt
    },
    basicblocklabel/.style={
        font=\normalsize,
        text=black!70
    },
    arrow/.style={
        -{Latex[length=1.8mm]},
        thick,
        draw=black!75
    },
    shortcut/.style={
        -{Latex[length=1.4mm]},
        dashed,
        thick,
        draw=purple!65!black
    },
    shortcutlabel/.style={
        font=\normalsize,
        text=purple!65!black,
        align=center
    },
    node distance=0.40cm
]

% Input and stem
\node[head] (input) {\rotatebox{90}{Image}};
\node[stem, right=of input] (conv1) {\rotatebox{90}{$7{\times}7$, 64, /2}};
\node[pool, right=of conv1] (pool1) {\rotatebox{90}{MaxPool, /2}};

% Layer 1: two BasicBlocks
\node[lone, right=1.10cm of pool1] (l10c1) {\rotatebox{90}{$3{\times}3$, 64}};
\node[lone, right=of l10c1]        (l10c2) {\rotatebox{90}{$3{\times}3$, 64}};
\node[lone, right=1.10cm of l10c2] (l11c1) {\rotatebox{90}{$3{\times}3$, 64}};
\node[lone, right=of l11c1]        (l11c2) {\rotatebox{90}{$3{\times}3$, 64}};
% Layer 2
\node[ltwo, right=1.60cm of l11c2] (l20c1) {\rotatebox{90}{$3{\times}3$, 128, /2}};
\node[ltwo, right=of l20c1]        (l20c2) {\rotatebox{90}{$3{\times}3$, 128}};
\node[ltwo, right=1.10cm of l20c2] (l21c1) {\rotatebox{90}{$3{\times}3$, 128}};
\node[ltwo, right=of l21c1]        (l21c2) {\rotatebox{90}{$3{\times}3$, 128}};
% Layer 3
\node[lthree, right=1.60cm of l21c2] (l30c1) {\rotatebox{90}{$3{\times}3$, 256, /2}};
\node[lthree, right=of l30c1]         (l30c2) {\rotatebox{90}{$3{\times}3$, 256}};
\node[lthree, right=1.10cm of l30c2]  (l31c1) {\rotatebox{90}{$3{\times}3$, 256}};
\node[lthree, right=of l31c1]         (l31c2) {\rotatebox{90}{$3{\times}3$, 256}};
% Layer 4
\node[lfour, right=1.60cm of l31c2] (l40c1) {\rotatebox{90}{$3{\times}3$, 512, /2}};
\node[lfour, right=of l40c1]         (l40c2) {\rotatebox{90}{$3{\times}3$, 512}};
\node[lfour, right=1.10cm of l40c2]  (l41c1) {\rotatebox{90}{$3{\times}3$, 512}};
\node[lfour, right=of l41c1]         (l41c2) {\rotatebox{90}{$3{\times}3$, 512}};
% Head
\node[head, right=1.60cm of l41c2] (avg) {\rotatebox{90}{AvgPool}};
\node[head, right=of avg]           (fc)  {\rotatebox{90}{FC}};

% BasicBlock containers
\begin{pgfonlayer}{background}

\node[basicblockbox, fit=(l10c1)(l10c2), label={[basicblocklabel]below:Block 1}] (bb10) {};
\node[basicblockbox, fit=(l11c1)(l11c2), label={[basicblocklabel]below:Block 2}] (bb11) {};

\node[basicblockbox, fit=(l20c1)(l20c2), label={[basicblocklabel]below:Block 1}] (bb20) {};
\node[basicblockbox, fit=(l21c1)(l21c2), label={[basicblocklabel]below:Block 2}] (bb21) {};

\node[basicblockbox, fit=(l30c1)(l30c2), label={[basicblocklabel]below:Block 1}] (bb30) {};
\node[basicblockbox, fit=(l31c1)(l31c2), label={[basicblocklabel]below:Block 2}] (bb31) {};

\node[basicblockbox, fit=(l40c1)(l40c2), label={[basicblocklabel]below:Block 1}] (bb40) {};
\node[basicblockbox, fit=(l41c1)(l41c2), label={[basicblocklabel]below:Block 2}] (bb41) {};

\end{pgfonlayer}

% Main arrows
\draw[arrow] (input)  -- (conv1);
\draw[arrow] (conv1)  -- (pool1);

\draw[arrow] (pool1)  -- (l10c1);
\draw[arrow] (l10c1)  -- (l10c2);
\draw[arrow] (l10c2)  -- (l11c1);
\draw[arrow] (l11c1)  -- (l11c2);

\draw[arrow] (l11c2)  -- (l20c1);
\draw[arrow] (l20c1)  -- (l20c2);
\draw[arrow] (l20c2)  -- (l21c1);
\draw[arrow] (l21c1)  -- (l21c2);

\draw[arrow] (l21c2)  -- (l30c1);
\draw[arrow] (l30c1)  -- (l30c2);
\draw[arrow] (l30c2)  -- (l31c1);
\draw[arrow] (l31c1)  -- (l31c2);

\draw[arrow] (l31c2)  -- (l40c1);
\draw[arrow] (l40c1)  -- (l40c2);
\draw[arrow] (l40c2)  -- (l41c1);
\draw[arrow] (l41c1)  -- (l41c2);

\draw[arrow] (l41c2)  -- (avg);
\draw[arrow] (avg)    -- (fc);

% Projection shortcuts used by the downsampling block in layers 2--4.
\draw[shortcut]
    ($(bb20.south west)+(0.18,0.28)$) --
    node[shortcutlabel, above=0.03cm, pos=0.50] {$1{\times}1$ proj, /2}
    ($(bb20.south east)+(-0.18,0.28)$);
\draw[shortcut]
    ($(bb30.south west)+(0.18,0.28)$) --
    node[shortcutlabel, above=0.03cm, pos=0.50] {$1{\times}1$ proj, /2}
    ($(bb30.south east)+(-0.18,0.28)$);
\draw[shortcut]
    ($(bb40.south west)+(0.18,0.28)$) --
    node[shortcutlabel, above=0.03cm, pos=0.50] {$1{\times}1$ proj, /2}
    ($(bb40.south east)+(-0.18,0.28)$);

% Stage braces and labels
\draw[decorate, decoration={brace, amplitude=5pt}]
    ($(bb10.north west)+(0,0.25)$) -- ($(bb11.north east)+(0,0.25)$)
    node[midway, yshift=0.48cm] {Layer 1};

\draw[decorate, decoration={brace, amplitude=5pt}]
    ($(bb20.north west)+(0,0.25)$) -- ($(bb21.north east)+(0,0.25)$)
    node[midway, yshift=0.48cm] {Layer 2};

\draw[decorate, decoration={brace, amplitude=5pt}]
    ($(bb30.north west)+(0,0.25)$) -- ($(bb31.north east)+(0,0.25)$)
    node[midway, yshift=0.48cm] {Layer 3};

\draw[decorate, decoration={brace, amplitude=5pt}]
    ($(bb40.north west)+(0,0.25)$) -- ($(bb41.north east)+(0,0.25)$)
    node[midway, yshift=0.48cm] {Layer 4};

\end{tikzpicture}%
%
}
\caption{Detailed ResNet-18 structure with BasicBlock containers. Each residual stage contains two BasicBlocks, and each BasicBlock contains two convolution operations. The first block of layers 2--4 includes a projection shortcut for downsampling.}
\label{fig:resnet18_detailed_basicblock_containers}
\end{figure}

%% RQ1.4 Kubernetes 5-service pod/service placement figure
\begin{figure}[htbp]
\centering
\resizebox{0.98\textwidth}{!}{%
\begin{tikzpicture}[
    font=\small,
    cluster/.style={
        draw=black!55,
        fill=black!2,
        rounded corners=8pt,
        thick,
        inner xsep=18pt,
        inner ysep=20pt
    },
    nodebox/.style={
        draw=black!65,
        fill=black!4,
        rounded corners=7pt,
        thick,
        inner xsep=16pt,
        inner ysep=18pt
    },
    pod/.style={
        draw=blue!65!black,
        fill=blue!10,
        rounded corners=4pt,
        minimum width=2.15cm,
        minimum height=1.05cm,
        align=center,
        thick
    },
    clientpod/.style={
        draw=orange!75!black,
        fill=orange!13,
        rounded corners=4pt,
        minimum width=2.25cm,
        minimum height=1.05cm,
        align=center,
        thick
    },
    svc/.style={
        draw=green!45!black,
        fill=green!12,
        rounded corners=3pt,
        minimum width=1.75cm,
        minimum height=0.55cm,
        align=center,
        thick,
        font=\scriptsize
    },
    arrow/.style={
        -{Latex[length=2mm]},
        thick,
        draw=black!70
    },
    returnarrow/.style={
        -{Latex[length=2mm]},
        thick,
        draw=black!55,
        dashed
    },
    smalllabel/.style={
        font=\scriptsize,
        text=black!65,
        align=center
    },
    segmentlabel/.style={
        font=\scriptsize,
        text=blue!65!black,
        align=center
    },
    transferlabel/.style={
        font=\scriptsize,
        text=red!70!black,
        align=center
    },
    node distance=0.55cm
]

% Top row: client and inference pods
\node[clientpod] (client) {Benchmark\\client pod};

\node[pod, right=1.05cm of client] (p1) {Pod 1\\\texttt{stem+layer1}};
\node[pod, right=0.95cm of p1] (p2) {Pod 2\\\texttt{layer2}};
\node[pod, right=0.95cm of p2] (p3) {Pod 3\\\texttt{layer3}};
\node[pod, right=0.95cm of p3] (p4) {Pod 4\\\texttt{layer4}};
\node[pod, right=0.95cm of p4] (p5) {Pod 5\\\texttt{avgpool+fc}};

% Services below each inference pod
\node[svc, below=0.45cm of p1] (s1) {Service 1};
\node[svc, below=0.45cm of p2] (s2) {Service 2};
\node[svc, below=0.45cm of p3] (s3) {Service 3};
\node[svc, below=0.45cm of p4] (s4) {Service 4};
\node[svc, below=0.45cm of p5] (s5) {Service 5};

% Background boxes: one Kubernetes node inside one cluster
\begin{pgfonlayer}{background}
\node[
    nodebox,
    fit=(client)(p1)(p2)(p3)(p4)(p5)(s1)(s2)(s3)(s4)(s5),
    label={[font=\small\bfseries, text=black!70]above:Single Kubernetes node}
] (k8snode) {};

\node[
    cluster,
    fit=(k8snode),
    label={[font=\small\bfseries, text=black!70]above:Local Kubernetes cluster / namespace \texttt{rq14}}
] (clusterbox) {};
\end{pgfonlayer}

% Request and forwarding path
\draw[arrow] (client) -- node[smalllabel, above] {gRPC} (p1);
\draw[arrow] (p1) -- node[transferlabel, above] {784 KiB} (p2);
\draw[arrow] (p2) -- node[transferlabel, above] {392 KiB} (p3);
\draw[arrow] (p3) -- node[transferlabel, above] {196 KiB} (p4);
\draw[arrow] (p4) -- node[transferlabel, above] {98 KiB} (p5);

% Pod-to-service visual association
\draw[black!35, thick] (s1) -- (p1);
\draw[black!35, thick] (s2) -- (p2);
\draw[black!35, thick] (s3) -- (p3);
\draw[black!35, thick] (s4) -- (p4);
\draw[black!35, thick] (s5) -- (p5);

% Return path
\draw[returnarrow]
    (p5.south) .. controls +(0,-1.25) and +(0,-1.25) ..
    node[smalllabel, below] {final output returned to caller}
    (client.south);

% Explanatory labels
\node[
    smalllabel,
    below=0.55cm of s3,
    text width=9.5cm
] {
    \texttt{chain\_5svc} is the maximum coarse-stage decomposition:
    each residual-stage boundary becomes a synchronous gRPC hop between pods.
};

\node[
    segmentlabel,
    above=0.35cm of p3,
    text width=7.5cm
] {
    Forward path: one request enters the first pod and intermediate activations are forwarded synchronously until inference completes.
};

\end{tikzpicture}%%
}
\caption{Kubernetes pod and service placement for the RQ1.4 five-service chain. The benchmark client and all inference pods are placed on the same Kubernetes node, while each model segment is exposed through a Kubernetes Service and communicates using synchronous gRPC forwarding. The labels above the arrows show the activation-transfer size at each coarse ResNet-18 boundary.}
\label{fig:rq14_chain5_kubernetes_placement}
\end{figure}